% =====================================================================
%  5. Step-size analysis and stability
% =====================================================================
\section{How big a step can we take?}

\begin{frame}{Start with quadratics, where everything is exact}
  Take $f(\x) = \tfrac12\x\T\bA\x - \vb\T\x$ with $\bA\succ0$ symmetric, so $\nabla f = \bA\x - \vb$ and $\bA\x^\star = \vb$. Then
  \[
    \e_{k+1} = \x_{k+1} - \x^\star = (\bI - \alpha\bA)\,\e_k \;\Longrightarrow\; \e_k = (\bI-\alpha\bA)^k\e_0 .
  \]
  Diagonalise $\bA = \mathbf Q\boldsymbol\Lambda\mathbf Q\T$ and write $\mathbf z_k = \mathbf Q\T\e_k$. The coupled iteration
  \alert{falls apart into $n$ independent scalar ones}:
  \[
    z_i^{(k)} = (1 - \alpha\lambda_i)^k\, z_i^{(0)}, \qquad i = 1,\dots,n .
  \]
  \begin{keybox}[Why bother with quadratics?]
    Close to a strict minimiser, $f(\x)\approx f(\x^\star) + \tfrac12(\x-\x^\star)\T\nabla^2 f(\x^\star)(\x-\x^\star)$.
    So the quadratic analysis with $\bA = \nabla^2 f(\x^\star)$ predicts how GD behaves \emph{near the end}
    on any smooth $f$ (Problem M1(d) uses exactly this).
  \end{keybox}
\end{frame}

\begin{frame}{Theorem: the stable window and the best fixed step}
  \begin{columns}[T]
    \begin{column}{0.52\textwidth}
      \begin{thmbox}[Theorem]
        \small
        GD converges from every $\x_0$ iff $0 < \alpha < 2/\lmax$, at rate
        \[
          \rho(\alpha) = \max\bigl\{|1-\alpha\lmin|,\ |1-\alpha\lmax|\bigr\}.
        \]
        The best choice is $\alpha^\star = \frac{2}{\lmin+\lmax}$, giving $\rho^\star = \frac{\kappa-1}{\kappa+1}$ with $\kappa = \frac{\lmax}{\lmin}$.
      \end{thmbox}
      \footnotesize\textbf{Proof.} $|1-\alpha\lambda|$ is convex in $\lambda$, so its maximum over $[\lmin,\lmax]$ sits at an end.
      Both ends are below 1 iff $\alpha<2/\lmax$. The best $\alpha$ makes the two ends equal: $1-\alpha\lmin = \alpha\lmax - 1$. $\blacksquare$
    \end{column}
    \begin{column}{0.46\textwidth}
      \centering
      \includegraphics[width=\linewidth]{fig_amplification}

      \vspace{2pt}
      \footnotesize To reach accuracy $\varepsilon$ you need about $\frac{\kappa+1}{2}\ln\frac1\varepsilon$ steps, i.e.\ $O\!\bigl(\kappa\log\tfrac1\varepsilon\bigr)$.
    \end{column}
  \end{columns}
\end{frame}

\begin{frame}{Same problem, four step sizes}
  \begin{columns}[T]
    \begin{column}{0.6\textwidth}
      \centering
      \includegraphics[width=\linewidth,height=0.80\textheight,keepaspectratio]{fig_stepsize_compare}
    \end{column}
    \begin{column}{0.38\textwidth}
      \small
      $\lambda(\bA)\in\{2,4\}$, so GD is stable only for $\alpha<0.5$.
      \begin{itemize}\footnotesize
        \item $\alpha=0.05$: $\rho = 0.9$. It gets there, eventually.
        \item $\alpha = \tfrac13$: $\rho = \tfrac13$, the M3 optimum.
        \item $\alpha = 0.48$: $|1-4\alpha| = 0.92$, rattling across the valley.
        \item $\alpha = 0.52$: $|1-4\alpha| = 1.08$, and the swings grow without bound.
      \end{itemize}
      \footnotesize\textcolor{mGrey}{Too small wastes time; too large wrecks the run.}
    \end{column}
  \end{columns}
\end{frame}

\begin{frame}{Seen from numerical ODEs: Euler's stability region}
  \begin{columns}[T]
    \begin{column}{0.5\textwidth}
      \small
      For a quadratic the gradient flow is $\dot\e = -\bA\e$, so each mode is the Dahlquist test equation
      $\dot y = \lambda y$ with $\lambda = -\lambda_i < 0$.
      \begin{itemize}\footnotesize
        \item The \textbf{exact flow} decays for every $\lambda_i>0$. Nothing can go wrong in continuous time.
        \item \textbf{Forward Euler} is stable only when $z = h\lambda$ lies in the disk $|1+z|<1$.
              Here $z_i = -\alpha\lambda_i$, so we need $0<\alpha\lambda_i<2$, which is exactly $\alpha<2/\lmax$.
        \item This is \textbf{stiffness}: the \emph{fastest} mode ($\lmax$) limits the step, while the \emph{slowest}
              ($\lmin$) decides how long you wait. Any explicit scheme pays roughly $\kappa$ steps.
      \end{itemize}
    \end{column}
    \begin{column}{0.48\textwidth}
      \centering
      % Explicit-Euler stability disk |1 + z| < 1 with the points z_i = -alpha*lambda_i
% (Problem M3: lambda = 2, 4) for alpha = 1/3 (stable) and alpha = 0.6 (unstable)
\begin{tikzpicture}[scale=1.35, >={Stealth[length=5pt]}]
  \shade[inner color=mGreen!28, outer color=mGreen!6] (-1,0) circle (1);
  \draw[line width=1.1pt, mGreen] (-1,0) circle (1);
  % axes
  \draw[->, mGrey] (-2.85,0) -- (0.6,0) node[right, font=\scriptsize] {$\mathrm{Re}\,z$};
  \draw[->, mGrey] (0,-1.2) -- (0,1.25) node[above, font=\scriptsize] {$\mathrm{Im}\,z$};
  \foreach \t/\lab in {-2/-2, -1/-1} {
    \draw[mGrey] (\t,0.05) -- (\t,-0.05);
    \node[font=\tiny, mGrey, fill=white, inner sep=1pt] at (\t,-0.17) {$\lab$};
  }
  \node[font=\tiny, mGrey] at (0.1,-0.15) {$0$};
  \node[font=\scriptsize\bfseries, mGreen] at (-1,0.55) {stable};
  % alpha = 1/3: z = -2/3, -4/3
  \fill[mBlue, draw=white, line width=0.6pt] (-0.667,0) circle (2.3pt);
  \fill[mBlue, draw=white, line width=0.6pt] (-1.333,0) circle (2.3pt);
  % alpha = 0.6: z = -1.2, -2.4
  \draw[mRed, line width=1.1pt, fill=white] (-1.2,0) circle (2.6pt);
  \draw[mRed, line width=1.1pt, fill=white] (-2.4,0) circle (2.6pt);
  \draw[mRed, ->] (-2.4,-0.62) node[below, font=\tiny] {escapes} -- (-2.4,-0.12);
\end{tikzpicture}


      \footnotesize
      \textcolor{mBlue}{$\bullet$} $\alpha = \frac13$: $z = -\frac23, -\frac43$, both inside.\quad
      \textcolor{mRed}{$\circ$} $\alpha = 0.6$: $z = -1.2, -2.4$, one escapes.

      \vspace{2pt}
      Backward Euler, $\x_{k+1} = \x_k - h\nabla f(\x_{k+1})$, is stable for every $h$.
      That is the proximal-point method; the catch is that each step is itself an optimisation problem.
    \end{column}
  \end{columns}
\end{frame}
